\subsection{Vanishing of the degree-\(163\) conjugation obstruction}
\label{sec:163-concrete-odd-inertia}

Put
\[
 R=R_{\rm global},\qquad \Gamma=G/R,\qquad
 P=P_{163},\qquad A=Q_{163}P,\qquad U=(P/A)^\vee.
\]
The module \(U\) is the nontrivial irreducible two-dimensional
\(\mathbb F_2\)-module for the cyclotomic quotient
\[
 \Gamma\xrightarrow{\lambda\bmod163}C_{162}\longrightarrow C_3;
\]
see Lemma~\ref{lem:163-integral-pair-decomposition}.
The omitted-prime theorem applies to the exact defining
intersection of \(R\), independently of its full degree-\(163\) image.

\begin{theorem}\label{thm:163-concrete-odd-inertia}
For every \(j\ge0\),
\begin{equation}
                         H^j_{\rm cont}(\Gamma,U)=0.
 \label{eq:163-concrete-cohomology-vanishing}
\end{equation}
\end{theorem}

\begin{proof}
Theorem~\ref{thm:simultaneous-global-cm-normalization} gives
\[
 R=R_{\rm quad}\cap
       \bigcap_{q\in\Pi_{\rm global}}\ker h_q,\qquad
 G=RH,\qquad 163\notin\Pi_{\rm global}.
\]
A primitive generator \(s\) of \(C_{162}\) acts on \(U\)
through a nontrivial element of \(C_3\).
Its minimal polynomial is \(X^2+X+1\), so \(s-1\) is
bijective on \(U\).
Theorem~\ref{thm:omitted-prime-vanishing}, with
\(\Lambda=\Pi_{\rm global}\), \(p=163\), and \(W=U\),
therefore proves
\eqref{eq:163-concrete-cohomology-vanishing}.
\end{proof}

\begin{corollary}\label{cor:163-single-extension-obstruction}
Suppose \(h_{163}(R)=P\), and put
\(J=\ker(h_{163}:R\to P)\).
For \(\kappa\in\mathscr X(J)\), the second obstruction
\(o_1(\kappa)\) of
Proposition~\ref{prop:163-two-stage-extension} is always
zero whenever it is defined.
The character \(\kappa\) extends to a continuous
character of \(G\) vanishing on \(H\) if and only if
\[
                   o_2(\kappa)=0
                       \quad\text{in }\mathbb F_2\,\mathfrak q.
\]
When it exists, this extension is unique.
Equivalently, restriction identifies
\begin{equation}
 \operatorname{Hom}_{\rm cts}(B,\mathbb F_2)
   \simeq
 \ker\bigl(o_2:\mathscr X(J)\longrightarrow
                               \mathbb F_2\,\mathfrak q\bigr).
 \label{eq:163-remaining-relative-obstruction}
\end{equation}
\end{corollary}

\begin{proof}
Theorem~\ref{thm:163-concrete-odd-inertia} gives
\(H^1(\Gamma,U)=0\), so the conjugation obstruction in
Proposition~\ref{prop:163-two-stage-extension} vanishes.
That proposition supplies the remaining condition,
uniqueness, and injectivity of restriction.
Characters of \(G\) vanishing on \(H\) are the
continuous characters of the relative abelian quotient
\(B\), giving \eqref{eq:163-remaining-relative-obstruction}.
\end{proof}

The theorem is independent of which full-image
alternative in
Proposition~\ref{prop:cm163-full-image-dichotomy} holds.
It does not decide between \(h_{163}(R)=A\) and
\(h_{163}(R)=P\), and does not show that the relative
quaternion obstruction \(o_2\) vanishes.
Thus the original question concerning
\(\operatorname{Hom}_{\rm cts}(B,\mathbb F_2)\) remains open.
